%! Simplifying Radicals; Adding \& Subtracting — Unit 5, Lesson 5.2 — Lesson Plan
\documentclass[10pt]{article}
\usepackage{algebra2-article}
\usepackage{algebra2-boxes}
\usepackage{graphicx}   % -article does NOT load graphicx; needed for thumbnails/screenshots
\graphicspath{{images/}}
\newcommand{\TallMath}[1]{$\displaystyle #1\rule[-1.4em]{0pt}{3.2em}$}
\newcommand{\CourseName}{Algebra 2: Shepherd}
\newcommand{\SchoolYear}{2026--2027}
\newcommand{\MeetingLength}{55 minutes}
\newcommand{\UnitNumberName}{Unit 5: Radical Functions \quad}
\newcommand{\LessonNumberName}{Lesson 5.2: Simplifying Radicals; Adding \& Subtracting}

\begin{document}
\begin{center}
    {\Large\bfseries \CourseName: \SchoolYear \\
    {\small\bfseries \UnitNumberName \LessonNumberName}}
\end{center}

\begin{tcolorbox}[colback=forestbg, colframe=forest, boxrule=0.9pt, arc=2mm,
    left=2mm, right=2mm, top=1.5mm, bottom=1.5mm]
    \textbf{Primary Objective:} Students can state the \emph{product} and \emph{quotient} properties
    of radicals and derive both from the exponent rules of Lesson 5.1, rather than memorizing them;
    can explain --- and prove --- why there is \emph{no} corresponding rule for addition; can write
    any radical in \emph{simplest form}, with numeric radicands, algebraic radicands, and indices
    above $2$, by pulling out the \emph{largest} perfect $n$th power; can simplify a variable
    radicand by dividing the exponent by the index, and explain that move as rewriting an improper
    fraction as a \emph{mixed number}; and can \emph{add and subtract} radicals by collecting
    \emph{like radicals} --- simplifying every term \emph{before} deciding whether any two of them
    are alike.
    \\[2pt]
    \textbf{Standards (2023 VA SOL):} \textbf{A2.EO.2a} --- simplify and determine equivalent
    radical expressions that include numeric and algebraic radicands --- is the governing standard
    and is assessed in every component. \textbf{A2.EO.2b} is served in its \emph{additive} half
    (add and subtract radical expressions, simplifying the result); multiplication, division, and
    rationalizing are deliberately held for Lesson 5.3, which is why this pair of lessons splits
    A2.EO.2b down the middle. \textbf{A2.EO.2c} (radical $\leftrightarrow$ rational exponent) is
    revisited as the \emph{source} of both properties and again in the mixed-number explanation of
    the exponent-division rule. The Activity's Part 3 and the homework's pendulum model reach
    \textbf{A2.F.2f} (evaluate a function and interpret it in context) and preview
    \textbf{A2.F.1b/c} (vertical stretch of a radical parent) ahead of Lesson 5.4.
\end{tcolorbox}

\renewcommand{\arraystretch}{1.4}
\begin{skillbox}[Priority Ideas \& Skills]{goldbox}
\begin{tabularx}{\linewidth}{@{}X|X@{}}
\textbf{Priority Ideas \& Skills} & \textbf{Key Understandings (the ``why'')} \\[2pt]
\hline
\vspace{-6pt}
\begin{itemize}[leftmargin=1.1em, nosep, after=\vspace{2pt}]
  \item Apply the \textbf{product} and \textbf{quotient} properties of radicals in both directions.
  \item Write a numeric radical in \textbf{simplest form} by extracting the \textbf{largest} perfect $n$th power.
  \item Simplify with \textbf{indices above $2$}, grouping factors in $n$'s rather than pairs.
  \item Simplify \textbf{algebraic} radicands by dividing the exponent by the index (quotient out, remainder in).
  \item Recognize \textbf{like radicals} --- same index \emph{and} same radicand --- and collect them.
  \item \textbf{Simplify first, decide second}: never declare ``cannot combine'' before simplifying.
  \item Reject $\sqrt{a+b}=\sqrt a+\sqrt b$, and \emph{prove} it fails for all positive $a,b$.
\end{itemize}
&
\vspace{-6pt}
\begin{itemize}[leftmargin=1.1em, nosep, after=\vspace{2pt}]
  \item These are not new rules. $\sqrt[n]{ab}=(ab)^{1/n}=a^{1/n}b^{1/n}$ --- the product property \emph{is} yesterday's power-of-a-product rule.
  \item A radical splits over \textbf{multiplication} and never over \textbf{addition}, and the arithmetic settles it in five seconds: $\sqrt9+\sqrt{16}=7$ while $\sqrt{9+16}=5$.
  \item The \textbf{index sets the group size}. $\sqrt[3]{9}$ is finished; a perfect square is worth nothing under a cube root.
  \item Dividing the exponent by the index \emph{is} converting $\frac{7}{2}$ to $3\frac12$ --- so ``quotient out, remainder in'' is a fact about mixed numbers, not a rule to memorize.
  \item Radicals combine exactly like terms: the radical plays the role of the variable, and only the coefficients move.
  \item Two radicals can be like radicals \textbf{in disguise} --- which is precisely why simplifying and combining belong in one lesson.
\end{itemize}
\end{tabularx}
\end{skillbox}

\begin{skillbox}[{Vocabulary, Concepts, \& Theorems}]{sky}
\renewcommand{\arraystretch}{1.3}
\begin{tabularx}{\linewidth}{@{}l X@{}}
\textbf{Product Property} & \TallMath{\sqrt[n]{ab}=\sqrt[n]{a}\cdot\sqrt[n]{b}} --- follows from $(ab)^{1/n}=a^{1/n}b^{1/n}$. For an \emph{even} index, $a,b\ge0$ (Lesson 5.1). \\
\textbf{Quotient Property} & \TallMath{\sqrt[n]{\tfrac{a}{b}}=\tfrac{\sqrt[n]{a}}{\sqrt[n]{b}}}, $b\ne0$ --- take the root of the top and of the bottom. \\
\textbf{No sum property} & $\sqrt{a+b}\ne\sqrt{a}+\sqrt{b}$. Squaring the false equation leaves $2\sqrt{ab}=0$, so it holds \emph{only} if $a=0$ or $b=0$. \\
\textbf{Simplest radical form} & (1) no factor of the radicand is a perfect $n$th power; (2) no fraction under the radical; (3) no radical in a denominator. Condition 3 is Lesson 5.3's job. \\
\textbf{Index sets the group size} & Factors escape a square root in pairs, a cube root in threes, a fourth root in fours. So $\sqrt[3]{9}$ is already simplest. \\
\textbf{Exponent-division rule} & $\sqrt[n]{x^{m}}$: divide $m$ by $n$ --- the quotient is the exponent outside, the remainder the exponent inside. Equivalently, $x^{m/n}$ as a \emph{mixed number}. \\
\textbf{Like radicals} & Same \emph{index} and same \emph{radicand}. Only like radicals combine, and only the coefficients change: $3\sqrt5+7\sqrt5=10\sqrt5$, never $10\sqrt{10}$. \\
\textbf{Rationalizing} & Clearing a radical out of a denominator --- named today, taught in Lesson 5.3. \\
\end{tabularx}
\end{skillbox}

\begin{fixedskillbox}[Activate Prior Knowledge \& Spiral Review]{forestbg}
    The Warm-Up rehearses three prerequisites, and each is a piece of today's lesson in disguise.
    \textbf{(1)} Perfect powers read backwards, then the \emph{largest perfect-square factor} of
    $72$, $50$, $48$ --- literally step one of the simplifying algorithm, extracted so that the
    algorithm itself has nothing new in it. Push on \emph{largest}: a student who writes
    $72=4\cdot18$ is not wrong but will finish at $2\sqrt{18}$ and believe they are done.
    \textbf{(2)} Combining like terms ($3x+7x$, $3x+7y$) with the follow-up question ``what exactly
    had to match?'' --- get the words \emph{the variable part} said aloud now, because
    \emph{the radical part} is then a two-word edit. \textbf{(3)} The exponent rules: verifying
    $(4\cdot9)^{1/2}=4^{1/2}\cdot9^{1/2}$ numerically \emph{before} the product property is stated,
    plus writing $\frac72$ as a mixed number and simplifying $x^{3}\cdot x^{1/2}$ --- which is the
    surprise ending of Section 4 pre-assembled. Note who is slow on the mixed number; they will
    struggle with algebraic radicands, not with the arithmetic ones.
\end{fixedskillbox}

\begin{skillbox}[Hook]{forestbg}
    Put three statements on the board and hand out calculators --- this is the one day in the unit
    when decimal approximation is the right tool, because the whole point is that a numeric check
    settles the question instantly. \textbf{``Two of these are true. One is a lie. Find it.''}
    \begin{center}
    (a) \; $\sqrt{8}\cdot\sqrt{2}=4$ \qquad (b) \; $\sqrt{9}+\sqrt{16}=\sqrt{25}$ \qquad
    (c) \; $\sqrt{8}+\sqrt{18}=5\sqrt{2}$
    \end{center}
    Give them three minutes. Most will convict (c) on sight --- it looks like nonsense --- and will
    be startled when $2.828+4.243=7.071$ and $5\sqrt2=7.071$ agree to three decimals. Meanwhile (b)
    looks reasonable and is the lie: $3+4=7$, not $5$. Then ask the question the whole lesson hangs
    on: \textbf{``What do the two \emph{true} ones have in common that the false one doesn't?''}
    Land it: the true ones involve \textbf{multiplication} inside the radical --- $8\cdot2$
    directly, and $8=4\cdot2$, $18=9\cdot2$ underneath. \emph{A radical splits over multiplication
    and never over addition.} Then flag the strange consequence in (c), which is why today's two
    skills are taught in one period: $\sqrt8$ and $\sqrt{18}$ share nothing visible, and they are
    the same radical in disguise. You cannot tell until you simplify. Yesterday the index told them
    what a radical \emph{is}; today it tells them how far a radical can be taken apart.
\end{skillbox}

\begin{skillbox}[Lesson]{forestbg}
\begin{multicols}{2}
\textbf{1. Where the properties come from.} Do not present them as new. Write
$\sqrt[n]{ab}=(ab)^{1/n}=a^{1/n}b^{1/n}=\sqrt[n]{a}\sqrt[n]{b}$ and let students see that the only
tool used is Warm-Up item 3. State the even-index fine print explicitly --- $a,b\ge0$, or the
pieces are not real --- and connect it to the 5.1 exit ticket: this is where a student who cannot
tell when a root exists will cheerfully turn $\sqrt{-8}$ into $\sqrt{-4}\cdot\sqrt2$. Then close
the Hook loop with the two-line table: $\sqrt{9\cdot16}=12=\sqrt9\cdot\sqrt{16}$, but
$\sqrt{9+16}=5\ne7=\sqrt9+\sqrt{16}$.

\textbf{2. Simplest form.} Give the three conditions, and say out loud that condition 3 is
tomorrow's job --- students should know something is being deferred rather than discover a rule
they were never taught. Work $\sqrt{72}=6\sqrt2$ two ways: largest factor ($36\cdot2$) and the
``I only saw $4\cdot18$'' route, which lands on the same answer through $2\sqrt{18}$. Name the
habit: \emph{the last step of every problem is checking the leftover radicand for another perfect
power}. Offer prime factorization as the safety net.

\columnbreak

\textbf{3. Higher indices.} One sentence carries this: \textbf{the index sets the group size}.
Pairs escape a square root, triples escape a cube root. Do $\sqrt[3]{54}$, $\sqrt[3]{16}$,
$\sqrt[4]{48}$, and $\sqrt[3]{-24}$ (odd index, negative radicand --- a 5.1 callback worth naming).
Then the trap: $\sqrt[3]{9}$ is \emph{already} simplified, even though $9$ is a perfect square.
Expect the majority to get this wrong on first contact.

\textbf{4. Algebraic radicands.} Teach the procedure (divide the exponent by the index; quotient
out, remainder in), then spend the real time on \emph{why}: $\sqrt{x^7}=x^{7/2}=x^{3\frac12}
=x^{3}\cdot x^{1/2}=x^{3}\sqrt x$. \emph{Simplifying a radical and rewriting an improper fraction
as a mixed number are the same move.} This is the lesson's intellectual center; students who see it
stop guessing at which power comes out. Finish with $\sqrt{50x^5y^8}$ and $\sqrt[3]{40a^7b^3}$,
handling the number and each variable separately. Note once that positive variables are assumed
(5.1's $|a|$ rule is set aside, not repealed).

\textbf{5. Quotient property --- keep it short.} Six quick items, two of them run
\emph{backwards} ($\sqrt{50}/\sqrt2=5$). Flag $\sqrt{3/5}$ as unfinished business and move on.

\textbf{6. Like radicals.} Return to Warm-Up item 2 and build the parallel column by column:
$3x+7x=10x$ beside $3\sqrt5+7\sqrt5=10\sqrt5$; $3x+7y$ beside $3\sqrt5+7\sqrt2$. Insist that only
the coefficients move --- $10\sqrt5$, never $10\sqrt{10}$. Then the Hook's payoff:
$\sqrt8+\sqrt{18}=2\sqrt2+3\sqrt2=5\sqrt2$, and the rule that governs the rest of the unit ---
\textbf{simplify first, decide second}.
\end{multicols}
\end{skillbox}

\begin{skillbox}[Active Monitoring]{redbox}
Circulate during Guided Practice and the tiers. Watch for and correct:
\begin{itemize}[leftmargin=1.2em, nosep]
  \item \textbf{stopping early} --- $\sqrt{72}=2\sqrt{18}$, $\sqrt{80}=2\sqrt{20}$; not wrong, just unfinished, and the fix is a habit (``check the leftover''), not a re-teach;
  \item \textbf{splitting over addition} --- $\sqrt{20}+\sqrt5=\sqrt{25}$; refute it with the student's own calculator every single time, because $6.71\ne5$ is a better argument than any sentence;
  \item \textbf{pairs under a cube root} --- $\sqrt[3]{54}=\sqrt[3]{9\cdot6}=3\sqrt[3]{6}$, which \emph{looks} right because a $3$ appears outside; ask where that $3$ came from;
  \item \textbf{adding the radicands too} --- $3\sqrt7+2\sqrt7=5\sqrt{14}$;
  \item \textbf{deciding too early} --- writing ``cannot combine'' for $\sqrt8+\sqrt{18}$ or $\sqrt{18}+\sqrt{32}$;
  \item \textbf{over-applying 5.0} --- marking $\sqrt[3]{-135}$ undefined because ``negatives are not allowed'' (odd index; they are);
  \item on algebraic radicands, \textbf{taking the whole exponent out} --- $\sqrt{x^{7}}=x^{7}$ or $x^{3.5}$ instead of $x^{3}\sqrt x$.
\end{itemize}
Cold-call prompts: ``Is there still a perfect square in there?'' ``What is the index, and how many
equal factors do you need to get one out?'' ``Check that on your calculator --- do both sides
agree?'' ``Are these like radicals? How do you know \emph{yet}?'' ``What is $7$ divided by $2$ as a
mixed number, and where do those two pieces go?''
\end{skillbox}

\begin{skillbox}[Group Work \& Differentiation]{redbox}
\textbf{Hiding in Plain Sight} activity (tiered; mirrors the notes).
\begin{multicols}{3}
\textbf{Tier R --- Remediate.} A six-row scaffold (radical $\to$ largest perfect-square factor $\to$
rewrite $\to$ simplified) over $\sqrt{32},\sqrt{44},\sqrt{75},\sqrt{90},\sqrt{128}$; a
\emph{finished-or-not} sort whose key item is $\sqrt[3]{9}$ (done --- a perfect square buys nothing
under a cube root); a ``find the families'' task where six radicals must be simplified before they
can be sorted into a $\sqrt3$ group and a $\sqrt2$ group and summed; and one concept question on
why $10\sqrt3+11\sqrt2$ stops there.

\columnbreak
\textbf{Tier A --- Approaching.} Six full simplifications including $\sqrt[3]{81y^{8}}$ and two
quotient-property items; four sums requiring simplification first; a three-row \textbf{error
analysis} (split over addition, added radicands, perfect square under a cube root) --- the highest
value items on the page, worth a whole-class share; and an exact-versus-approximate perimeter
problem ($\sqrt{75}$ by $\sqrt{12}$ meters $\Rightarrow 14\sqrt3\approx24.2$ m) that ends by asking
why a builder wants the exact form.

\columnbreak
\textbf{Tier E --- Extension.} Three parts. \textbf{(1)} A genuine proof that
$\sqrt a+\sqrt b=\sqrt{a+b}$ forces $ab=0$ --- square, cancel, and the surviving $2\sqrt{ab}$ is
exactly what the classic error deletes. \textbf{(2)} The staircase
$\sqrt2+\sqrt8+\sqrt{18}+\sqrt{32}+\sqrt{50}=15\sqrt2$, generalized to
$\frac{n(n+1)}{2}\sqrt2$ after noticing $2k^{2}$, then rebuilt over $3$. \textbf{(3)} A bridge to
5.4: a table showing $f(x)=\sqrt{8x}$ is always twice $g(x)=\sqrt{2x}$, proved by
$\sqrt{8x}=2\sqrt{2x}$ --- simplifying a radical reveals a \textbf{vertical stretch}, with the
domain and the endpoint unchanged.
\end{multicols}
\end{skillbox}

\begin{skillbox}[Individual Work \& Assessment]{redbox}
\textbf{Exit Ticket} (independent, no notes, no calculator): four simplifications spanning a
numeric radical, a cube root, an algebraic radicand, and a sum that only combines after both terms
are simplified ($2\sqrt{12}+\sqrt{27}=7\sqrt3$); an \textbf{explanation} item requiring the student
to refute $\sqrt9+\sqrt{16}=\sqrt{25}$ \emph{numerically} and then state the property that does
hold; and an \textbf{SOL-style multiple-choice} item on $\sqrt{50}+\sqrt{32}$ (answer $9\sqrt2$)
whose distractors each break exactly one idea --- $\sqrt{82}$ (added the radicands), $9\sqrt4$
(simplified correctly, then added the radicands anyway), $20\sqrt2$ (multiplied the coefficients).
Collect and sort into ``got it / stopped early / adders / did not simplify first.'' The last two
piles govern how 5.3 opens: distributing over radicals magnifies both errors immediately, since
$\sqrt2\left(\sqrt6+\sqrt{10}\right)$ requires simplifying \emph{after} multiplying, not before.
\end{skillbox}

\begin{skillbox}[Reinforcement \& Extension]{goldbox}
\textbf{Homework:} eight numeric simplifications spanning indices $2$ through $5$ and including a
negative radicand under an odd index; six algebraic radicands, one with a fraction; eight
add-or-subtract items, two of which are answered ``cannot combine'' \emph{only after} being
simplified (and one, $\sqrt8+\sqrt[3]8$, which quietly rewards noticing that $\sqrt[3]8$ is just
$2$); three quotient-property items; a two-part conceptual item --- refute
$\sqrt{36+64}=\sqrt{36}+\sqrt{64}$ numerically, and explain why $\sqrt[3]{16}$ is unfinished while
$\sqrt[3]{9}$ is done; and a closing item that simplifies $\sqrt{x^{9}}$ \emph{twice}, once by
splitting the radicand and once via $x^{9/2}=x^{4\frac12}$, asking why the two roads always agree.
\textbf{Extension:} the pendulum, $T(L)=2\pi\sqrt{L/32}$, with lengths $2,8,18,32$ ft chosen so the
fraction reduces to a perfect square and the periods come out exactly $\frac\pi2,\pi,\frac{3\pi}2,
2\pi$. Its shape deliberately matches 5.1's Kepler model: a table of exact values, a growth-rate
reading (quadruple the length to double the period --- what a $\frac12$ exponent does), a
clockmaker who wrongly plans to double the length, and a structural rewrite
$\sqrt{L/32}=\sqrt{2L}/8$ obtained by making the denominator a perfect square --- rationalizing,
performed before the word exists. \textbf{Preview:} Lesson 5.3, \emph{Multiplying \& Dividing
Radicals; Rationalizing} --- the product property run \emph{forwards}, distributing and FOIL over
radicals, and finally simplest-form condition 3, cleared with \textbf{conjugates} (the same move
that made $a+bi$ real back in Lesson 2.4).
\end{skillbox}
\end{document}
